\documentclass[twocolumn]{aastex631}
\usepackage{amsmath}
\usepackage{booktabs}
\shorttitle{SDSS target vector for feedback-model validation}
\shortauthors{Simulation validation}
\begin{document}

\title{SDSS target vector for feedback-model validation: selection-aware SDSS optical proxy integration}
\author{NebulaMind Research Autopilot}
\affiliation{Public SDSS DR17 data only}

\begin{abstract}
We use the cached SDSS DR17 emission-line subset to define a compact optical target vector for forward-model validation. The pilot records quenched fraction, optical AGN incidence, and color versus mass/redshift across cells with sufficient counts, providing an observed target vector for later simulation work. It remains an empirical denominator study rather than a direct simulation comparison.
\end{abstract}

\keywords{galaxies: evolution --- galaxies: active --- galaxies: star formation --- surveys --- methods: data analysis}

\section{Purpose and claim contract}\label{sec:purpose}
This draft preserves the active proposal title, 'Forward-modelled validation of cosmological feedback prescriptions', but narrows the supported claim to the cached SDSS optical measurement named in the results. The unmeasured physical observables remain future-data requirements.

The claim contract is intentionally conservative. Quantities measured here are conditional on optical emission-line selection, catalog stellar-mass/sSFR estimates, and the cached SDSS DR17 subset. Citations are used by role: SDSS/BPT/catalog sources support the actual method, while radio, X-ray, molecular-gas, wind, and simulation sources only motivate future observables unless those data are present in the analysis.


\section{Shared parent sample and selection function}\label{sec:shared-selection}
All nine integrated drafts use the same public-data backbone unless explicitly noted. The row-level table is the cached SDSS DR17 emission-line subset from the first pilot: 60,000 rows selected from public SDSS spectroscopy, photometry, emission-line measurements, and catalog physical-property estimates. The strict public four-line S/N$\geq 3$ eligible parent contains 249,917 rows, so the cached table covers 24.0\% of that strict parent. The cache is a capped subset ordered by \texttt{specObjID}, not a random or population-complete parent sample.

\begin{deluxetable*}{lrrr}
\tabletypesize{\scriptsize}
\tablecaption{Shared SDSS DR17 selection cascade used before paper-specific quantities.\label{tab:selection-cascade}}
\tablehead{\colhead{Selection stage} & \colhead{Public DR17 rows} & \colhead{Cached rows} & \colhead{Retention vs. spectro-z parent}}
\startdata
SpecObj GALAXY, 0.02<z<0.12 & 501,060 & -- & 1.000 \\
plus galSpecInfo/PhotoObj/galSpecExtra and mass/sSFR bounds & 416,554 & -- & 0.831 \\
plus galSpecLine join & 416,554 & -- & 0.831 \\
four BPT lines positive with positive errors & 373,445 & 60,000 & 0.745 \\
four BPT lines S/N>=3 & 249,917 & 60,000 & 0.499 \\
four BPT lines S/N>=5 & 176,523 & 42,446 & 0.352 \\
four BPT lines S/N>=10 & 91,768 & 22,311 & 0.183 \\
\enddata
\tablecomments{Counts are read-only public SDSS DR17 count queries plus the cached local CSV. Cached rows are shown only where the cache applies.}
\end{deluxetable*}

The four-line requirement is strongly selection dependent. In the public counts, S/N$\geq3$ keeps 33.6\% of the $-12<\log {\rm sSFR}<-11$ parent bin but 94.9\% of the $-10<\log {\rm sSFR}<-9.5$ bin. Therefore every incidence, quenching, density, gas-denominator, or target-vector statement in this integration is conditional on the four-line emission-line selection.

Cached-versus-public marginal checks found no redshift, stellar-mass, or sSFR bin with a cached-minus-public fraction difference above 5 percentage points. The largest absolute differences were 2.03 percentage points in redshift, -1.63 percentage points in stellar mass, and -0.58 percentage points in sSFR. This is a representativeness diagnostic only; it does not make the capped cache random or complete.


\section{Measurements}\label{sec:measurements}
The row-level measurements include redshift, stellar mass, catalog specific star-formation rate, model colors, and four optical line fluxes/errors. BPT labels and all derived denominators are recomputed locally from the cached table. Catalog SFR/sSFR values are treated as low-redshift SDSS physical-property estimates rather than direct resolved gas or feedback measurements \citep{sdssdr17,brinchmann2004,york2000}.


\section{Optical target vector for simulation validation}\label{sec:topic-result}
The consolidated proposal question is: What compact SDSS target vector of quenched fraction, optical AGN incidence, and colour versus mass/redshift can be used for forward-model validation? The integrated manuscript deliberately demotes the claim to the directly measured SDSS quantity. It is not a full physical-feedback test.

\begin{itemize}
\item The pilot writes 15 mass-redshift cells with n >= 50 as a compact validation vector.
\item Across mass bins, quenched fractions span 0.005-0.729; optical AGN fractions span 0.003-0.520.
\item The output is an observed target vector for simulation forward modelling, not a direct simulation comparison.
\end{itemize}


\begin{figure}
\centering
\includegraphics[width=\columnwidth]{../figures/fig-topic.pdf}
\caption{SDSS DR17 optical denominator/proxy diagnostic for the feedback-model validation target vector. The figure summarizes the cached optical result used for target definition.}
\label{fig:topic}
\end{figure}

\section{Interpretation and missing observables}\label{sec:missing}
SDSS-only pilot; full proposal requires the additional survey data named in the research-topic page. The full proposal requires: simulation mocks passed through the SDSS/MaNGA/ALMA/X-ray/radio selection functions and aperture/noise models.

Simulation suites and mock-observation methods define the future comparison problem; no simulation mock has been forward-modelled or ranked in this pilot \citep{tng2019,eagle2015,simba2019,imanga2023,donnari2021,dubois2013,dubois2016}.


\section{Reproducibility and safety}\label{sec:repro}
This manuscript uses the cached SDSS DR17 subset and the accompanying local manifest of inputs and outputs. The output is a local draft PDF and manifest entry only. No public-linked PDF was replaced.

\section{Conclusion}\label{sec:conclusion}
The pilot writes 15 mass-redshift cells with $n \geq 50$ as a compact validation vector, with quenched fractions spanning 0.005--0.729 and optical AGN fractions spanning 0.003--0.520. This observed target vector is useful for simulation forward modelling, but it still requires mock-observation pipelines before any model comparison can be claimed.


\begin{thebibliography}{}
\bibitem[Abdurro'uf et al.(2022)]{sdssdr17} Abdurro'uf, Accetta, K., Aerts, C., et al. 2022, ApJS, 259, 35
\bibitem[Brinchmann et al.(2004)]{brinchmann2004} Brinchmann, J., Charlot, S., White, S.~D.~M., et al. 2004, MNRAS, 351, 1151
\bibitem[York et al.(2000)]{york2000} York, D.~G., Adelman, J., Anderson, J.~E., Jr., et al. 2000, AJ, 120, 1579

\bibitem[Dave et al.(2019)]{simba2019} Dave, R., Angles-Alcazar, D., Narayanan, D., et al. 2019, MNRAS, 486, 2827
\bibitem[Donnari et al.(2021)]{donnari2021} Donnari, M., Pillepich, A., Nelson, D., et al. 2021, MNRAS, 506, 4760
\bibitem[Dubois et al.(2013)]{dubois2013} Dubois, Y., Gavazzi, R., Peirani, S., \& Silk, J. 2013, MNRAS, 433, 3297
\bibitem[Dubois et al.(2016)]{dubois2016} Dubois, Y., Peirani, S., Pichon, C., et al. 2016, MNRAS, 463, 3948
\bibitem[Nanni et al.(2023)]{imanga2023} Nanni, L., Thomas, D., Trayford, J., et al. 2023, MNRAS, 518, 2605
\bibitem[Nelson et al.(2019)]{tng2019} Nelson, D., Springel, V., Pillepich, A., et al. 2019, Computational Astrophysics and Cosmology, 6, 2
\bibitem[Schaye et al.(2015)]{eagle2015} Schaye, J., Crain, R.~A., Bower, R.~G., et al. 2015, MNRAS, 446, 521
\end{thebibliography}

\end{document}
